\chapter{Categories}

En aquest capítol resolem amb detall una selecció d'exercicis corresponents al capítol \emph{Categories} de la part de Teoria. L'objectiu no és només arribar al resultat, sinó mostrar \emph{com} es raona en el llenguatge categòric: comprovar axiomes, seguir camins, traduir diagrames a equacions i reconèixer una mateixa estructura en conjunts, preordres, lògica i grafs. Convé intentar cada exercici abans de llegir-ne la solució.

\section{Definició de categoria}

\begin{exercici}\label{exr:pr-rel}
Comprova que els conjunts amb les relacions formen una categoria $\cat{Rel}$, on un morfisme $A\to B$ és una relació $R\subseteq A\times B$, la identitat de $A$ és la relació d'igualtat, i la composició de $R\subseteq A\times B$ amb $S\subseteq B\times C$ és
\[
S\comp R=\{(a,c)\in A\times C\mid \exists\,b\in B\text{ tal que }(a,b)\in R\text{ i }(b,c)\in S\}.
\]
\end{exercici}

\begin{solucio}
Hem de comprovar els dos axiomes. Comencem per les \textbf{identitats}. Sigui $\id_A=\{(a,a)\mid a\in A\}$ la relació d'igualtat sobre $A$. Per a tota relació $R\subseteq A\times B$,
\[
R\comp\id_A
=\{(a,b)\mid \exists\,a'\ (a,a')\in\id_A\text{ i }(a',b)\in R\}
=R,
\]
perquè $(a,a')\in\id_A$ obliga $a'=a$. Anàlogament $\id_B\comp R=R$.

Per a l'\textbf{associativitat}, siguin $R\subseteq A\times B$, $S\subseteq B\times C$ i $T\subseteq C\times D$. Un parell $(a,d)$ pertany a $T\comp(S\comp R)$ si i només si existeixen $b\in B$ i $c\in C$ tals que
\[
(a,b)\in R,
\qquad
(b,c)\in S,
\qquad
(c,d)\in T.
\]
La mateixa condició caracteritza $(T\comp S)\comp R$. Per tant
\[
T\comp(S\comp R)=(T\comp S)\comp R,
\]
i $\cat{Rel}$ és una categoria.
\end{solucio}

\section{Fletxes, camins i diagrames commutatius}

\begin{exercici}
Llegeix el diagrama següent i escriu l'equació que expressa que commuta:
\[
\begin{tikzpicture}[baseline=(m.center)]
  \matrix (m) [matrix of math nodes, row sep=2.7em, column sep=3.6em] {
    A & B \\
    C & D \\
  };
  \draw[-{Stealth[scale=1.1]}] (m-1-1) -- node[above] {$f$} (m-1-2);
  \draw[-{Stealth[scale=1.1]}] (m-1-1) -- node[left] {$u$} (m-2-1);
  \draw[-{Stealth[scale=1.1]}] (m-1-2) -- node[right] {$v$} (m-2-2);
  \draw[-{Stealth[scale=1.1]}] (m-2-1) -- node[below] {$g$} (m-2-2);
\end{tikzpicture}
\]
Explica també per què l'ordre de les composicions és el que escrius.
\end{exercici}

\begin{solucio}
Hi ha dos camins dirigits de $A$ a $D$. El camí superior-dret és
\[
A\xrightarrow{f}B\xrightarrow{v}D,
\]
i per tant determina $v\comp f$. El camí esquerre-inferior és
\[
A\xrightarrow{u}C\xrightarrow{g}D,
\]
i determina $g\comp u$. El quadrat commuta exactament quan
\[
\boxed{v\comp f=g\comp u}.
\]
L'ordre de cada composició es llegeix de dreta a esquerra: en $v\comp f$ seguim primer $f$ i després $v$.
\end{solucio}

\begin{exercici}
Considera a $\cat{Set}$ quatre còpies de $\mathbb Z$ i les aplicacions
\[
f(n)=n+1,
\qquad
u(n)=2n,
\qquad
v(n)=2n,
\qquad
g(n)=n+2.
\]
Comprova que el quadrat
\[
\begin{tikzpicture}[baseline=(m.center)]
  \matrix (m) [matrix of math nodes, row sep=2.8em, column sep=4.2em] {
    \mathbb Z & \mathbb Z \\
    \mathbb Z & \mathbb Z \\
  };
  \draw[-{Stealth[scale=1.05]}] (m-1-1) -- node[above] {$f$} (m-1-2);
  \draw[-{Stealth[scale=1.05]}] (m-1-1) -- node[left] {$u$} (m-2-1);
  \draw[-{Stealth[scale=1.05]}] (m-1-2) -- node[right] {$v$} (m-2-2);
  \draw[-{Stealth[scale=1.05]}] (m-2-1) -- node[below] {$g$} (m-2-2);
\end{tikzpicture}
\]
commuta.
\end{exercici}

\begin{solucio}
Cal comprovar que $v\comp f=g\comp u$. Per a qualsevol $n\in\mathbb Z$,
\[
(v\comp f)(n)=v(n+1)=2n+2,
\]
mentre que
\[
(g\comp u)(n)=g(2n)=2n+2.
\]
Les dues aplicacions coincideixen en tots els enters. Per tant
\[
v\comp f=g\comp u,
\]
i el quadrat commuta. Aquest exemple mostra com una afirmació sobre funcions es pot llegir simultàniament com una igualtat algebraica i com la commutativitat d'un diagrama.
\end{solucio}

\section{Primers exemples}

\begin{exercici}
Sigui $T$ un sistema deductiu amb una relació de deduïbilitat $\vdash_T$ reflexiva i transitiva. Comprova amb detall que les fórmules, amb una única fletxa $A\to B$ quan $A\vdash_T B$, formen una categoria prima $\cat{Prov}_T$.
\end{exercici}

\begin{solucio}
Els \textbf{objectes} són les fórmules. Entre dues fórmules $A$ i $B$ posem un únic morfisme $A\to B$ exactament quan $A\vdash_T B$.

La \textbf{identitat} existeix perquè la deduïbilitat és reflexiva:
\[
A\vdash_T A.
\]
Per tant hi ha una fletxa $\id_A\colon A\to A$.

La \textbf{composició} existeix perquè la deduïbilitat és transitiva. Si
\[
A\vdash_T B
\qquad\text{i}\qquad
B\vdash_T C,
\]
aleshores
\[
A\vdash_T C.
\]
Diagramàticament,
\[
\begin{tikzpicture}[baseline=(m.center)]
  \matrix (m) [matrix of math nodes, column sep=3.4em] { A & B & C \\ };
  \draw[-{Stealth[scale=1.05]}] (m-1-1) -- (m-1-2);
  \draw[-{Stealth[scale=1.05]}] (m-1-2) -- (m-1-3);
  \draw[-{Stealth[scale=1.0]}] (m-1-1) to[bend right=25] (m-1-3);
\end{tikzpicture}
\]

Finalment, entre dues fórmules hi ha com a màxim una fletxa. Això fa que les igualtats exigides pels axiomes d'associativitat i identitat siguin automàtiques sempre que les composicions existeixin. Per tant $\cat{Prov}_T$ és una categoria prima.

Cal remarcar què hem oblidat: si hi ha dues demostracions diferents de $B$ a partir d'$A$, aquesta categoria les identifica en una sola fletxa. Només conserva la relació de deduïbilitat.
\end{solucio}

\begin{exercici}
Siguin els grafs dirigits $G$ i $H$ següents:
\[
G:\quad 0\xrightarrow{a}1,
\qquad
H:\quad
\begin{tikzpicture}[baseline=(m.center)]
  \matrix (m) [matrix of math nodes, row sep=2.0em, column sep=2.8em] {
    & y \\
    x & z \\
  };
  \draw[-{Stealth[scale=1.0]}] (m-2-1) -- node[above left] {$r$} (m-1-2);
  \draw[-{Stealth[scale=1.0]}] (m-2-1) -- node[below] {$s$} (m-2-2);
\end{tikzpicture}
\]
Defineix $F_V(0)=x$, $F_V(1)=y$ i $F_E(a)=r$. Comprova que això defineix un morfisme de grafs $F\colon G\to H$. Explica per què substituir $F_E(a)=r$ per $F_E(a)=s$, mantenint $F_V$, no defineix un morfisme de grafs.
\end{exercici}

\begin{solucio}
En $G$ tenim
\[
s_G(a)=0,
\qquad
t_G(a)=1,
\]
i en $H$,
\[
s_H(r)=x,
\qquad
t_H(r)=y.
\]
Per tant
\[
F_V(s_G(a))=F_V(0)=x=s_H(r)=s_H(F_E(a))
\]
i
\[
F_V(t_G(a))=F_V(1)=y=t_H(r)=t_H(F_E(a)).
\]
Així es compleixen les dues equacions
\[
F_V\comp s_G=s_H\comp F_E,
\qquad
F_V\comp t_G=t_H\comp F_E,
\]
i $F$ és un morfisme de grafs.

Si poséssim $F_E(a)=s$, l'origen encara seria correcte perquè $s_H(s)=x$, però el destí fallaria:
\[
t_H(s)=z\neq y=F_V(1).
\]
Per tant el quadrat corresponent al destí no commutaria i no tindríem un morfisme de grafs.
\end{solucio}

\section{Isomorfismes}

\begin{exercici}
En una categoria qualsevol, comprova que si $f\colon A\to B$ i $g\colon B\to C$ tenen inversos, aleshores $g\comp f$ també en té, i
\[
(g\comp f)^{-1}=f^{-1}\comp g^{-1}.
\]
\end{exercici}

\begin{solucio}
La situació és
\[
A\xrightarrow{f}B\xrightarrow{g}C.
\]
Per recórrer aquest camí en sentit contrari hem d'invertir també l'ordre: primer $g^{-1}$ i després $f^{-1}$. Proposem, doncs, $f^{-1}\comp g^{-1}$ com a invers.

D'una banda,
\[
(f^{-1}\comp g^{-1})\comp(g\comp f)
=f^{-1}\comp(g^{-1}\comp g)\comp f
=f^{-1}\comp f
=\id_A,
\]
i de l'altra,
\[
(g\comp f)\comp(f^{-1}\comp g^{-1})
=g\comp(f\comp f^{-1})\comp g^{-1}
=g\comp g^{-1}
=\id_C.
\]
Per tant $(g\comp f)^{-1}=f^{-1}\comp g^{-1}$.
\end{solucio}

\begin{exercici}
Comprova que a $\cat{Rel}$ una relació $R\colon A\to B$ és un isomorfisme si i només si és el graf d'una bijecció de $A$ a $B$.
\end{exercici}

\begin{solucio}
Si $f\colon A\to B$ és una bijecció i $R$ és el seu graf, el graf de $f^{-1}$ és una relació $B\to A$ i les dues composicions són les relacions d'igualtat. Per tant $R$ és un isomorfisme.

Recíprocament, suposem que $R$ té un invers $S\colon B\to A$, és a dir,
\[
S\comp R=\id_A,
\qquad
R\comp S=\id_B.
\]
Per a cada $a\in A$, com $(a,a)\in S\comp R$, existeix algun $b\in B$ amb $(a,b)\in R$ i $(b,a)\in S$. Així cada $a$ es relaciona amb almenys un $b$.

Si $(a,b)\in R$ i $(a,b')\in R$, triem $b_0$ amb $(a,b_0)\in R$ i $(b_0,a)\in S$. Llavors
\[
(b_0,b)\in R\comp S=\id_B,
\qquad
(b_0,b')\in R\comp S=\id_B,
\]
i per tant $b=b_0=b'$. Així cada $a$ es relaciona amb un únic $b$, i $R$ és el graf d'una aplicació $f\colon A\to B$.

Aplicat a $S$, el mateix argument dona una aplicació $g\colon B\to A$. Les igualtats $S\comp R=\id_A$ i $R\comp S=\id_B$ es converteixen en
\[
g\comp f=\id_A,
\qquad
f\comp g=\id_B.
\]
Per tant $f$ és bijectiva i $R$ n'és el graf.
\end{solucio}

\begin{exercici}
Sigui $\cat{C}$ una categoria. Comprova que la relació \enquote{ésser isomorf} entre objectes de $\cat{C}$ és una relació d'equivalència.
\end{exercici}

\begin{solucio}
\textbf{Reflexivitat}: $\id_A$ és el seu propi invers, de manera que $A\cong A$.

\textbf{Simetria}: si $f\colon A\to B$ és un isomorfisme amb invers $f^{-1}\colon B\to A$, aleshores $f^{-1}$ és també un isomorfisme, amb invers $f$; per tant $B\cong A$.

\textbf{Transitivitat}: si $f\colon A\to B$ i $g\colon B\to C$ són isomorfismes, la composició $g\comp f$ també ho és pel primer exercici de la secció; per tant $A\cong C$.
\end{solucio}

\begin{exercici}
A la categoria prima $\cat{Prov}_T$, comprova que dues fórmules $A$ i $B$ són isomorfes si i només si són interdeduïbles:
\[
A\vdash_T B
\qquad\text{i}\qquad
B\vdash_T A.
\]
\end{exercici}

\begin{solucio}
Si $A\cong B$, per definició existeixen fletxes
\[
A\longrightarrow B
\qquad\text{i}\qquad
B\longrightarrow A.
\]
A $\cat{Prov}_T$ això significa exactament $A\vdash_T B$ i $B\vdash_T A$.

Recíprocament, si es donen aquestes dues deduïbilitats, existeixen les dues fletxes. Com que $\cat{Prov}_T$ és prima, qualsevol endomorfisme $A\to A$ és necessàriament l'única fletxa $\id_A$, i igualment per a $B$. Per tant les dues composicions són les identitats i les fletxes són inverses. Així $A\cong B$.
\end{solucio}

\section{Classes especials de morfismes}

\begin{exercici}
Demostra que un morfisme $f$ és un isomorfisme si i només si és un monomorfisme escindit (té retracció) i un epimorfisme.
\end{exercici}

\begin{solucio}
Si $f$ és un isomorfisme, el seu invers és alhora retracció i secció, de manera que $f$ és mono escindit; i tot isomorfisme és epi. Recíprocament, suposem que $f\colon A\to B$ té una retracció $r$ (amb $r\comp f=\id_A$) i que és epi. Partint de $r\comp f=\id_A$ i component amb $f$ per l'esquerra,
\[
(f\comp r)\comp f = f\comp(r\comp f) = f\comp\id_A = f = \id_B\comp f.
\]
Com que $f$ és epi, es pot cancel·lar el factor $f$ de la dreta, d'on
\[
f\comp r=\id_B.
\]
Juntament amb $r\comp f=\id_A$, això mostra que $r$ és també invers per la dreta de $f$: $f$ és un isomorfisme amb invers $r=f^{-1}$.
\end{solucio}

\begin{exercici}
A $\cat{Set}$, considera $f\colon\{0,1\}\to\{\ast\}$, l'única aplicació, i $g\colon\{\ast\}\to\{0,1\}$ amb $g(\ast)=0$. Comprova quina de les dues composicions $f\comp g$ i $g\comp f$ és una identitat, digues quin dels dos morfismes és una secció i quin una retracció (i de qui), i determina quin dels dos és un monomorfisme i quin un epimorfisme.
\end{exercici}

\begin{solucio}
Calculem les dues composicions. Per una banda, $f\comp g\colon\{\ast\}\to\{\ast\}$ compleix $(f\comp g)(\ast)=f(0)=\ast$, de manera que $f\comp g=\id_{\{\ast\}}$. Per l'altra, $g\comp f\colon\{0,1\}\to\{0,1\}$ envia tant $0$ com $1$ a $0$, de manera que $g\comp f\neq\id_{\{0,1\}}$ (per exemple, $(g\comp f)(1)=0\neq 1$).

Així, l'única igualtat que val és $f\comp g=\id_{\{\ast\}}$. Llegida des de cada morfisme, aquesta equació diu dues coses alhora: $g$ és una \emph{secció} de $f$ (un invers per la dreta de $f$), i $f$ és una \emph{retracció} de $g$ (un invers per l'esquerra de $g$). Convé insistir que \enquote{secció} i \enquote{retracció} són papers relatius a l'altre morfisme, no propietats absolutes: aquí $g$ no és \enquote{una secció} en abstracte, sinó una secció \emph{de $f$}.

Pel que fa a mono i epi: com que $f$ té una secció, $f$ és un epimorfisme; com a aplicació, és exhaustiva, d'acord amb la caracterització dels epimorfismes de $\cat{Set}$. Com que $g$ té una retracció, $g$ és un monomorfisme; com a aplicació, és injectiva. En canvi, $f$ no és mono, perquè l'aplicació $f$ no és injectiva ($f(0)=f(1)$), i $g$ no és epi, perquè l'aplicació $g$ no és exhaustiva ($1\notin g(\{\ast\})$). Cap dels dos és un isomorfisme: ho serien només si $g\comp f$ fos també una identitat, cosa que hem vist que no passa.
\end{solucio}

\begin{exercici}\label{exr:pr-prima-mono-epi}
Sigui $\cat{C}$ una categoria prima, és a dir, una categoria en què entre dos objectes hi ha com a màxim un morfisme. Comprova que tot morfisme de $\cat{C}$ és alhora un monomorfisme i un epimorfisme. Explica per què aquest fet mostra que, en una categoria arbitrària, \enquote{mono} i \enquote{epi} no es poden interpretar com \enquote{injectiu} i \enquote{exhaustiu}.
\end{exercici}

\begin{solucio}
Sigui $f\colon A\to B$ un morfisme de $\cat{C}$. Per veure que és un \textbf{monomorfisme}, siguin $g,h\colon X\to A$ amb $f\comp g=f\comp h$:
\[
\begin{tikzpicture}[baseline=(m.center)]
  \matrix (m) [matrix of math nodes, column sep=3.4em] { X & A & B \\ };
  \draw[-{Stealth[scale=1.0]}] ([yshift=2.6pt]m-1-1.east) -- node[above] {$g$} ([yshift=2.6pt]m-1-2.west);
  \draw[-{Stealth[scale=1.0]}] ([yshift=-2.6pt]m-1-1.east) -- node[below] {$h$} ([yshift=-2.6pt]m-1-2.west);
  \draw[-{Stealth[scale=1.1]}] (m-1-2) -- node[above] {$f$} (m-1-3);
\end{tikzpicture}
\qquad
f\comp g=f\comp h\ \Longrightarrow\ g=h.
\]
Els morfismes $g$ i $h$ tenen el mateix domini $X$ i el mateix codomini $A$: tots dos pertanyen a $\Hom(X,A)$. Com que la categoria és prima, aquest homconjunt té com a màxim un element, i per tant $g=h$. Observem que ni tan sols hem fet servir la hipòtesi $f\comp g=f\comp h$: la cancel·lació es compleix perquè no hi ha res a cancel·lar.

Per veure que és un \textbf{epimorfisme}, siguin $g,h\colon B\to Y$ amb $g\comp f=h\comp f$:
\[
\begin{tikzpicture}[baseline=(m.center)]
  \matrix (m) [matrix of math nodes, column sep=3.4em] { A & B & Y \\ };
  \draw[-{Stealth[scale=1.1]}] (m-1-1) -- node[above] {$f$} (m-1-2);
  \draw[-{Stealth[scale=1.0]}] ([yshift=2.6pt]m-1-2.east) -- node[above] {$g$} ([yshift=2.6pt]m-1-3.west);
  \draw[-{Stealth[scale=1.0]}] ([yshift=-2.6pt]m-1-2.east) -- node[below] {$h$} ([yshift=-2.6pt]m-1-3.west);
\end{tikzpicture}
\qquad
g\comp f=h\comp f\ \Longrightarrow\ g=h.
\]
Ara $g,h\in\Hom(B,Y)$, i el mateix argument dona $g=h$. També ho podem obtenir per dualitat: la categoria oposada $\cat{C}\op$ també és prima, perquè $\Hom_{\cat{C}\op}(B,A)$ està en bijecció amb $\Hom_{\cat{C}}(A,B)$; per tant $f\op$ és mono a $\cat{C}\op$, és a dir, $f$ és epi a $\cat{C}$ (teorema~\ref{teo:dualitat}).

\emph{Per què això impedeix llegir mono com a injectiu i epi com a exhaustiu.} Hi ha dues raons, una més radical que l'altra.

Primer, en una categoria prima els morfismes no són, en general, aplicacions. En un preordre, la fletxa $x\to y$ només registra que $x\le y$; no envia cap element enlloc, i preguntar si és \enquote{injectiva} no té sentit. Mono i epi, en canvi, sempre en tenen, perquè es defineixen només amb la composició.

Segon, fins i tot si volguéssim mantenir l'analogia, entraria en contradicció. A $\cat{Set}$, una aplicació injectiva i exhaustiva és bijectiva, és a dir, un isomorfisme. Si \enquote{mono $=$ injectiu} i \enquote{epi $=$ exhaustiu} valguessin en general, tot morfisme d'una categoria prima seria un isomorfisme. Però a $\cat{2}=(0\to 1)$ la fletxa $u\colon 0\to 1$ és mono i epi pel que acabem de demostrar, i no té invers, perquè no hi ha cap fletxa $1\to 0$. En un preordre qualsevol passa el mateix: la fletxa $x\to y$ és sempre mono i epi, i només és un isomorfisme quan també $y\le x$.

La conclusió és que \enquote{mono $=$ injectiu} i \enquote{epi $=$ exhaustiu} són teoremes sobre $\cat{Set}$, no definicions. Les definicions són les de cancel·lació, i en una categoria prima es compleixen de la manera més trivial possible.
\end{solucio}

\begin{exercici}\label{pr-mono-no-escindit}
Compara les aplicacions injectives $i\colon\varnothing\to\{\ast\}$ i $j\colon\{0\}\hookrightarrow\{0,1\}$ de $\cat{Set}$. Totes dues són monomorfismes? Quina d'elles és un monomorfisme escindit?
\end{exercici}

\begin{solucio}
Totes dues són injectives i, per tant, monomorfismes de $\cat{Set}$.

La inclusió $j$ té una retracció: l'aplicació $r\colon\{0,1\}\to\{0\}$ constant compleix $r\comp j=\id_{\{0\}}$. És, doncs, un monomorfisme escindit.

L'aplicació $i$, en canvi, no en té cap: una retracció seria una aplicació $\{\ast\}\to\varnothing$, i no n'hi ha cap, perquè $\ast$ no té on anar. Per tant, $i$ és un monomorfisme que no és escindit. L'exemple mostra que l'afirmació \enquote{tota aplicació injectiva té una inversa per l'esquerra} només és certa per a dominis no buits; i aquest contraexemple no depèn de l'axioma de l'elecció, a diferència del cas dual de les seccions.
\end{solucio}

\section{Construccions sobre categories}

\begin{exercici}
Descriu explícitament la categoria producte $\cat{2}\times\cat{3}$, on $\cat{2}$ és la categoria associada a l'ordre $0<1$ i $\cat{3}$ la categoria associada a $0<1<2$. Quants objectes i quants morfismes té? Com es componen?
\end{exercici}

\begin{solucio}
Els objectes són les parelles $(i,j)$ amb $i\in\{0,1\}$ i $j\in\{0,1,2\}$. N'hi ha sis i es poden disposar en una graella:
\[
\begin{tikzpicture}[baseline=(m.center)]
  \matrix (m) [matrix of math nodes, row sep=2.4em, column sep=2.8em] {
    (0,0) & (0,1) & (0,2) \\
    (1,0) & (1,1) & (1,2) \\
  };
  \draw[-{Stealth[scale=1.05]}] (m-1-1) -- (m-1-2);
  \draw[-{Stealth[scale=1.05]}] (m-1-2) -- (m-1-3);
  \draw[-{Stealth[scale=1.05]}] (m-2-1) -- (m-2-2);
  \draw[-{Stealth[scale=1.05]}] (m-2-2) -- (m-2-3);
  \draw[-{Stealth[scale=1.05]}] (m-1-1) -- (m-2-1);
  \draw[-{Stealth[scale=1.05]}] (m-1-2) -- (m-2-2);
  \draw[-{Stealth[scale=1.05]}] (m-1-3) -- (m-2-3);
\end{tikzpicture}
\]
Només hi hem dibuixat les fletxes bàsiques; les composicions produeixen també fletxes més llargues i diagonals.

Com que $\cat{2}$ i $\cat{3}$ són primes, entre $(i,j)$ i $(i',j')$ hi ha un únic morfisme exactament quan
\[
i\le i'
\qquad\text{i}\qquad
j\le j'.
\]
$\cat{2}$ té tres morfismes i $\cat{3}$ en té sis. Com que un morfisme del producte és una parella, hi ha
\[
3\cdot6=18
\]
morfismes. La composició es fa component a component. Tots els quadrats de la graella commuten perquè el producte és també una categoria prima.
\end{solucio}

\begin{exercici}
Descriu explícitament la categoria oposada $\cat{P}\op$ quan $\cat{P}$ és un preordre $(P,\leq)$ vist com a categoria. Quina relació d'ordre li correspon?
\end{exercici}

\begin{solucio}
A $\cat{P}$ hi ha un morfisme $x\to y$ exactament quan $x\leq y$. A $\cat{P}\op$ les fletxes s'inverteixen:
\[
x\longrightarrow y\text{ a }\cat{P}\op
\quad\Longleftrightarrow\quad
y\longrightarrow x\text{ a }\cat{P}
\quad\Longleftrightarrow\quad
y\leq x.
\]
Per tant $\cat{P}\op$ és el preordre invers $(P,\geq)$. En aquest exemple, invertir les fletxes és literalment invertir la desigualtat.
\end{solucio}

\begin{exercici}
Sigui $\cat{2}=(0\to 1)$ la categoria amb dos objectes i una fletxa no trivial. Prenent com a ambient la mateixa $\cat{2}$, descriu explícitament la categoria sobre l'objecte $1$, $\cat{2}/1$, i la categoria sota l'objecte $1$, $1/\cat{2}$, comptant-ne objectes i morfismes. Comprova després que $(\cat{2}/1)\op$ \emph{no} és isomorfa a $1/\cat{2}$, i identifica amb quina categoria sí que ho és.
\end{exercici}

\begin{solucio}
Escrivim la fletxa no trivial de $\cat{2}$ com $u\colon 0\to 1$.

\emph{La categoria $\cat{2}/1$} (sobre $1$). Els seus objectes són els morfismes de $\cat{2}$ amb codomini $1$: la fletxa $u\colon 0\to 1$ i la identitat $\id_1\colon 1\to 1$. Són, doncs, \emph{dos} objectes. Un morfisme de $u$ a $\id_1$ és una fletxa $h$ de $\cat{2}$ amb $\id_1\comp h=u$, és a dir $h=u$; i les identitats de cada objecte completen els morfismes. La categoria $\cat{2}/1$ té la forma d'una fletxa entre dos objectes: és isomorfa a $\cat{2}$.

\emph{La categoria $1/\cat{2}$} (sota $1$). Els seus objectes són els morfismes de $\cat{2}$ amb domini $1$: només $\id_1$, perquè de $1$ no en surt cap altra fletxa. És, doncs, \emph{un} sol objecte, amb només la seva identitat: la categoria $\cat{1}$.

\emph{El contrast.} $\cat{2}/1$ té dos objectes i $1/\cat{2}$ en té un. Per tant $(\cat{2}/1)\op$ té dos objectes i \emph{no} pot ser isomorfa a $1/\cat{2}$. El que sí que val és la fórmula general de dualitat
\[
(\cat{2}/1)\op\;\cong\;1/(\cat{2}\op).
\]
En efecte, $\cat{2}\op$ és la fletxa invertida $1\to 0$; sota l'objecte $1$ a $\cat{2}\op$ hi ha els morfismes amb domini $1$, que són $\id_1$ i la fletxa $1\to 0$: \emph{dos} objectes, com $(\cat{2}/1)\op$. La moral: en dualitzar la construcció \enquote{sobre} en \enquote{sota} cal dualitzar també la categoria ambient, no deixar-la intacta.
\end{solucio}

\section{Principi de dualitat}

\begin{exercici}\label{exr:pr-dual-mono-epi}
Demostra que la composició de dos monomorfismes és un monomorfisme. A continuació, sense repetir el càlcul, dedueix per dualitat que la composició de dos epimorfismes és un epimorfisme. Indica exactament què s'inverteix en passar a $\cat{C}\op$ i què es manté.
\end{exercici}

\begin{solucio}
\emph{Primera part: el cas dels monomorfismes.} Siguin $f\colon A\to B$ i $g\colon B\to C$ monomorfismes, i siguin $u,v\colon X\to A$ amb $(g\comp f)\comp u=(g\comp f)\comp v$:
\[
\begin{tikzpicture}[baseline=(m.center)]
  \matrix (m) [matrix of math nodes, column sep=3em] { X & A & B & C \\ };
  \draw[-{Stealth[scale=1.0]}] ([yshift=2.6pt]m-1-1.east) -- node[above] {$u$} ([yshift=2.6pt]m-1-2.west);
  \draw[-{Stealth[scale=1.0]}] ([yshift=-2.6pt]m-1-1.east) -- node[below] {$v$} ([yshift=-2.6pt]m-1-2.west);
  \draw[-{Stealth[scale=1.1]}] (m-1-2) -- node[above] {$f$} (m-1-3);
  \draw[-{Stealth[scale=1.1]}] (m-1-3) -- node[above] {$g$} (m-1-4);
\end{tikzpicture}
\qquad
(g\comp f)\comp u=(g\comp f)\comp v\ \Longrightarrow\ u=v.
\]
Per associativitat, la hipòtesi és $g\comp(f\comp u)=g\comp(f\comp v)$. Com que $g$ és mono, en cancel·lem el factor de l'esquerra i obtenim $f\comp u=f\comp v$. Com que $f$ és mono, tornem a cancel·lar i obtenim $u=v$. Per tant $g\comp f$ és mono. Observem l'ordre: primer es cancel·la $g$, el factor \emph{exterior}, i després $f$.

\emph{Segona part: dualització.} Procedim en tres passos.

\textbf{1. L'enunciat és categòric.} \enquote{Si $f\colon A\to B$ i $g\colon B\to C$ són monomorfismes, aleshores $g\comp f$ és un monomorfisme} només parla d'objectes, morfismes, composició i igualtats entre composicions (la definició de mono és d'aquesta mena). Per tant té dual, i com que l'hem demostrat en una categoria qualsevol, el principi de dualitat (teorema~\ref{teo:dualitat}) garanteix que el dual també és vàlid en tota categoria.

\textbf{2. Què s'inverteix.} Recorrem l'enunciat element per element:
\begin{center}
\begin{tabular}{@{}ll@{}}
\toprule
A $\cat{C}$ & A l'enunciat dual \\
\midrule
$f\colon A\to B$, \quad $g\colon B\to C$ & $f\colon B\to A$, \quad $g\colon C\to B$ \\
la composta $g\comp f\colon A\to C$ & la composta $f\comp g\colon C\to A$ \\
fletxes de prova $u,v\colon X\to A$ & fletxes de prova $u,v\colon A\to X$ \\
mono: $m\comp u=m\comp v\Rightarrow u=v$ & epi: $u\comp m=v\comp m\Rightarrow u=v$ \\
\bottomrule
\end{tabular}
\end{center}
Es mantenen els quantificadors (\enquote{per a tot $u,v$}), la implicació, la igualtat i l'associativitat, que és autodual. L'enunciat dual diu, doncs: si $f\colon B\to A$ i $g\colon C\to B$ són epimorfismes, aleshores $f\comp g\colon C\to A$ és un epimorfisme. Rebatejant les fletxes perquè la cadena es llegeixi d'esquerra a dreta, $p\colon A\to B$ i $q\colon B\to C$, això és: \emph{si $p$ i $q$ són epimorfismes, $q\comp p$ és un epimorfisme.}

\textbf{3. Per què és vàlid.} És l'argument del principi, fet explícit. Siguin $p\colon A\to B$ i $q\colon B\to C$ epimorfismes de $\cat{C}$. A $\cat{C}\op$, les fletxes $p\op\colon B\to A$ i $q\op\colon C\to B$ són monomorfismes, perquè ser epi a $\cat{C}$ és ser mono a $\cat{C}\op$. Per la primera part, aplicada a $\cat{C}\op$, la composta $p\op\comp q\op\colon C\to A$ és mono a $\cat{C}\op$. Però $p\op\comp q\op=(q\comp p)\op$, de manera que $(q\comp p)\op$ és mono a $\cat{C}\op$, és a dir, $q\comp p$ és epi a $\cat{C}$.

\emph{Comprovació: la prova dual.} No l'hem necessitat, però val la pena escriure-la per veure que és la primera llegida a l'inrevés. Siguin $u,v\colon C\to X$ amb $u\comp(q\comp p)=v\comp(q\comp p)$:
\[
\begin{tikzpicture}[baseline=(m.center)]
  \matrix (m) [matrix of math nodes, column sep=3em] { A & B & C & X \\ };
  \draw[-{Stealth[scale=1.1]}] (m-1-1) -- node[above] {$p$} (m-1-2);
  \draw[-{Stealth[scale=1.1]}] (m-1-2) -- node[above] {$q$} (m-1-3);
  \draw[-{Stealth[scale=1.0]}] ([yshift=2.6pt]m-1-3.east) -- node[above] {$u$} ([yshift=2.6pt]m-1-4.west);
  \draw[-{Stealth[scale=1.0]}] ([yshift=-2.6pt]m-1-3.east) -- node[below] {$v$} ([yshift=-2.6pt]m-1-4.west);
\end{tikzpicture}
\qquad
u\comp(q\comp p)=v\comp(q\comp p)\ \Longrightarrow\ u=v.
\]
Per associativitat, $(u\comp q)\comp p=(v\comp q)\comp p$; com que $p$ és epi, $u\comp q=v\comp q$; com que $q$ és epi, $u=v$. Les mateixes passes, amb cada composició en l'ordre invers: ara es cancel·la primer $p$, el factor de la \emph{dreta}, que és la fletxa per on comença el camí.

\emph{Error típic.} En dualitzar és fàcil invertir les fletxes però oblidar d'invertir l'ordre de la composició, i escriure \enquote{$p\comp q$ és epi}. Si $p\colon A\to B$ i $q\colon B\to C$, aquesta composta ni tan sols està definida, tret que $C=A$. Comprovar els dominis i els codominis de cada composta és la millor manera de detectar l'error.
\end{solucio}

\section{Grafs i categories lliures}

\begin{exercici}
Considera el graf amb un sol vèrtex i dues arestes que són bucles, $a$ i $b$. Descriu la categoria lliure que genera. Quants morfismes té?
\end{exercici}

\begin{solucio}
La categoria lliure té un sol objecte. Els morfismes són tots els camins finits:
\[
\varepsilon,
\quad a,
\quad b,
\quad aa,
\quad ab,
\quad ba,
\quad bb,
\quad aaa,
\quad\ldots
\]
on $\varepsilon$ és el camí buit i actua com a identitat. La composició és la concatenació de paraules en ordre de recorregut: $b\comp a=ab$. Com que cap relació no identifica camins diferents, hi ha una infinitat numerable de morfismes. Llevat de l'isomorfisme canònic que inverteix l'ordre de les lletres (exemple~\ref{ex:monoide-lliure-cat}), és el monoide lliure sobre l'alfabet $\{a,b\}$.
\end{solucio}

\section{Qüestions de mida}

\begin{exercici}
En una categoria petita, l'oposada $\cat{C}\op$ també és petita. Comprova-ho, i comprova a més que $\cat{C}$ i $\cat{C}\op$ tenen exactament el mateix nombre de morfismes.
\end{exercici}

\begin{solucio}
Els objectes de $\cat{C}\op$ són els mateixos que els de $\cat{C}$. Cada morfisme $f\colon A\to B$ de $\cat{C}$ correspon al mateix morfisme vist com $f\colon B\to A$ a $\cat{C}\op$. Això dona una bijecció entre les col·leccions de morfismes de totes dues categories.

Per tant, si els objectes i els morfismes de $\cat{C}$ formen conjunts, també ho fan els de $\cat{C}\op$. En particular, $\cat{C}$ és petita si i només si ho és $\cat{C}\op$.
\end{solucio}

\begin{exercici}\label{exr:pr-set-localment-petita}
Comprova que $\cat{Set}$ és localment petita: per a conjunts $A$ i $B$, mostra que la col·lecció de les aplicacions $A\to B$ és un conjunt, identificant cada aplicació amb el seu graf, que és un subconjunt d'$A\times B$.
\end{exercici}

\begin{solucio}
Una aplicació $f\colon A\to B$ queda determinada pel seu \textbf{graf}
\[
\Gamma_f=\{(a,f(a))\mid a\in A\}\subseteq A\times B.
\]
Un subconjunt $\Gamma\subseteq A\times B$ és el graf d'alguna aplicació $A\to B$ exactament quan és \emph{funcional}: per a cada $a\in A$ hi ha un únic $b\in B$ amb $(a,b)\in\Gamma$. Així, $f\mapsto\Gamma_f$ és una bijecció entre les aplicacions $A\to B$ i els subconjunts funcionals d'$A\times B$; la inversa envia $\Gamma$ a l'aplicació que assigna a cada $a$ l'únic $b$ amb $(a,b)\in\Gamma$.

Vegem ara que aquests subconjunts funcionals formen un conjunt. Com que $A$ i $B$ són conjunts, el producte $A\times B$ és un conjunt, i també ho és el seu conjunt de parts $\mathcal P(A\times B)$. La condició \enquote{ser funcional} és una propietat ben definida dels elements de $\mathcal P(A\times B)$, de manera que, per l'axioma de separació,
\[
\{\Gamma\in\mathcal P(A\times B)\mid \forall a\in A\ \exists!\,b\in B\ (a,b)\in\Gamma\}
\]
és un conjunt. Per la bijecció anterior, $\Hom_{\cat{Set}}(A,B)$ és un conjunt.

Una precisió, la mateixa de l'observació~\ref{obs:extrems-implicits} de la Teoria per a $\cat{Rel}$: el graf no determina el codomini, perquè $\Gamma_f\subseteq A\times B$ també és un subconjunt d'$A\times B'$ per a tot $B'\supseteq B$. Per això, en rigor, un morfisme de $\cat{Set}$ és la terna $(A,B,\Gamma_f)$. Això no canvia la conclusió: les ternes $(A,B,\Gamma)$ amb $\Gamma$ funcional formen el conjunt $\{A\}\times\{B\}\times\{\Gamma\in\mathcal P(A\times B)\mid \Gamma\text{ funcional}\}$, que està en bijecció amb l'anterior.

\emph{Abast.} El mateix argument explica per què $\Hom(A,B)$ és un conjunt en els exemples habituals. A $\cat{Rel}$, directament, $\Hom_{\cat{Rel}}(X,Y)$ és, llevat dels extrems, $\mathcal P(X\times Y)$. A les categories de conjunts estructurats ($\cat{Grp}$, $\cat{Top}$, $\cat{Vect}_{\mathbb K}$, $\cat{Mon}$, $\cat{Pos}$), els morfismes són aplicacions entre els conjunts subjacents que compleixen una condició addicional; formen, doncs, un subconjunt de l'homconjunt de $\cat{Set}$ corresponent, de nou per separació. En canvi, la col·lecció de \emph{tots} els objectes de $\cat{Set}$ no és un conjunt (paradoxa de Russell): $\cat{Set}$ és localment petita sense ser petita.
\end{solucio}
